\chapter{Polynomials, Rings and Fields}
\label{ch:polynoms}

\medskip

\textbf{Notation.} $K[x]$ --- polynomials over field $K$; $\overline{K}$ --- algebraic closure.
$\deg f$ --- degree; $f'$ --- formal derivative $\bigl(f'=\sum i a_i x^{i-1}$, with the usual rules$\bigr)$.
$\operatorname{Res}(f,g)$ --- resultant; $\Delta(f)$ --- discriminant.
$e_k$ --- elementary symmetric sums; $p_k=\sum\alpha_i^k$ --- power sums.
$\mathbb{F}_q$ --- finite field of order $q=p^n$.
$[L:K]$ --- degree of field extension.
$\Phi_n(x)$ --- $n$-th cyclotomic polynomial.
A polynomial $f\in L[x]$ is \textit{separable} if it has no repeated roots; a field extension $L/K$ is separable if every algebraic element has a separable minimal polynomial.

\medskip

\section{Resultant and discriminant}
\label{sec:resultant-discriminant}

\subsection{Resultant --- definition and common root criterion}
\label{subsec:resultant-def}

\textit{Statement.} For $f(x)=a_nx^n+\dots+a_0,\;g(x)=b_mx^m+\dots+b_0$ over an integral domain: $\operatorname{Res}(f,g)=\det\operatorname{Syl}(f,g)$.

$$\operatorname{Res}(f,g)=0 \iff f,g\text{ have a common divisor of degree}\ge 1\text{ in }\overline{K}[x].$$

\textit{Idea.} The Sylvester matrix is $(m+n)\times(m+n)$: first $m$ rows contain shifted coefficients of $f$, next $n$ rows contain shifted coefficients of $g$:

\subsection{Resultant via roots}
\label{subsec:resultant-roots}

$f=a_n\prod(x-\alpha_i),\;g=b_m\prod(x-\beta_j)$:
$$\operatorname{Res}(f,g)=a_n^m b_m^n\prod_{i,j}(\alpha_i-\beta_j)=a_n^m\prod_i g(\alpha_i)=(-1)^{nm}b_m^n\prod_j f(\beta_j).$$

\textit{Properties.} Multiplicativity; $\operatorname{Res}(g,f)=(-1)^{nm}\operatorname{Res}(f,g)$; $\operatorname{Res}(f,x-a)=f(a)$.

\subsection{Discriminant}
\label{subsec:discriminant}

$\Delta(f)=\frac{(-1)^{n(n-1)/2}}{a_n}\operatorname{Res}(f,f')=a_n^{2n-2}\prod_{i<j}(\alpha_i-\alpha_j)^2.$

$\Delta(f)=0 \iff f$ has a multiple root. $\Delta(f(x+c))=\Delta(f)$, $\Delta(cf)=c^{2n-2}\Delta(f)$.

\subsection{Discriminant of $x^n+ax+b$}
\label{subsec:discriminant-trinomial}

$$\Delta(x^n+ax+b)=(-1)^{n(n-1)/2}\bigl((1-n)^{n-1}a^n+n^n b^{n-1}\bigr).$$
Check: $n=2$: $a^2-4b$; $n=3$: $-(4a^3+27b^2)$.

\subsection{Sign of the discriminant and number of real roots}
\label{subsec:discriminant-real}

$f\in\mathbb{R}[x]$ of degree $n$ with no multiple roots. Let $s$ be the number of pairs of complex conjugate roots. Then $\operatorname{sign}\Delta(f)=(-1)^s$.

\medskip

\section{Vieta's formulas, symmetric polynomials, Newton's identities}
\label{sec:vieta-symmetric}

\subsection{Vieta's formulas}
\label{subsec:vieta}

$f(x)=a_nx^n+\dots+a_0$ with roots $\alpha_1,\dots,\alpha_n$:
$$e_k(\alpha_1,\dots,\alpha_n)=(-1)^k\frac{a_{n-k}}{a_n}.$$
In particular, $\sum\alpha_i=-a_{n-1}/a_n$, $\prod\alpha_i=(-1)^n a_0/a_n$.

\subsection{Newton--Girard identities}
\label{subsec:newton-girard}

$p_k=\sum\alpha_i^k$. For $1\le k\le n$:
$$p_k=e_1 p_{k-1}-e_2 p_{k-2}+\dots+(-1)^{k-1}k\,e_k.$$
For $k>n$: $p_k=e_1 p_{k-1}-\dots+(-1)^{n-1}e_n p_{k-n}$.

\subsection{Fundamental theorem of symmetric polynomials}
\label{subsec:fundamental-symmetric}

Every symmetric $\Phi\in R[x_1,\dots,x_n]$ can be expressed uniquely in terms of $e_1,\dots,e_n$.

\subsection{Logarithmic derivative}
\label{subsec:log-derivative}

$\dfrac{f'(x)}{f(x)}=\sum_i\dfrac1{x-\alpha_i}$.
Consequence: $(n-1)(f')^2\ge n f f''$ when all roots are real.

\subsection{Newton and Maclaurin inequalities}
\label{subsec:newton-maclaurin}

If all roots of $f(x)=\prod(x+\alpha_i)$ are real, then $E_k=e_k/\binom{n}{k}$ satisfy $E_k^2\ge E_{k-1}E_{k+1}$ (log-concavity).
For $\alpha_i>0$: $S_1\ge S_2^{1/2}\ge\dots\ge S_n^{1/n}$, where $S_k=e_k/\binom{n}{k}$.

\subsection{Descartes' rule of signs}
\label{subsec:descartes}

The number of positive roots of $f\in\mathbb{R}[x]$ (counted with multiplicity) equals the number of sign changes in the coefficient sequence $V(f)$ or is less by an even number.

\subsection{Root bounds}
\label{subsec:root-bounds}

\begin{enumerate}
\item \textit{Cauchy:} all roots of $f$ satisfy $|z|\le 1+\max_{k<n}|a_k/a_n|$.
\item \textit{Lagrange:} all roots satisfy $|z|\le\max(1,\sum_{k=0}^{n-1}|a_k/a_n|)$.
\item \textit{Enestr\"om--Kakeya:} if $0<a_0\le a_1\le\dots\le a_n$, then all roots lie in $|z|\le 1$.
\end{enumerate}

\medskip

\section{Polynomials over $\mathbb{F}_p$ and $\mathbb{F}_q$}
\label{sec:polynomials-over-fq}

\subsection{Frobenius endomorphism}
\label{subsec:frobenius}

In characteristic $p$: $(x+y)^p=x^p+y^p$. The map $\varphi:a\mapsto a^p$ is a ring homomorphism. Over $\mathbb{F}_q$ ($q=p^n$) it is an automorphism generating $\operatorname{Gal}(\mathbb{F}_{p^n}/\mathbb{F}_p)$.

\subsection{Identity $x^q-x=\prod_{a\in\mathbb{F}_q}(x-a)$}
\label{subsec:xq-x}

$x^q-x=\prod_{a\in\mathbb{F}_q}(x-a)$, $x^{q-1}-1=\prod_{a\in\mathbb{F}_q^\times}(x-a)$. Consequence: $\gcd(f,x^q-x)=\prod_{a: f(a)=0}(x-a)$.

\subsection{Power sums $\sum_{a\in\mathbb{F}_q}a^k$}
\label{subsec:power-sums-fq}

$$\sum_{a\in\mathbb{F}_q}a^k=\begin{cases}-1,&(q-1)\mid k,\\0,&\text{otherwise}.\end{cases}$$

\subsection{``Too many roots'' principle}
\label{subsec:too-many-roots}

$f\in\mathbb{F}_q[x]$ of degree $n$ has at most $n$ roots.

\subsection{Cyclicity of $\mathbb{F}_q^\times$}
\label{subsec:fq-cyclic}

$\mathbb{F}_q^\times$ is cyclic of order $q-1$.

\subsection{Universal factorization of $x^{q^n}-x$}
\label{subsec:universal-factorization}

$$x^{q^n}-x=\prod_{d\mid n}\;\prod_{\substack{f\text{ monic irreducible}\\\deg f=d}} f(x).$$
The number of monic irreducibles of degree $n$ over $\mathbb{F}_q$:
$$N_q(n)=\frac1n\sum_{d\mid n}\mu(d)\,q^{n/d}>0.$$

\subsection{Chevalley--Warning}
\label{subsec:chevalley-warning}

$f_1,\dots,f_r\in\mathbb{F}_q[x_1,\dots,x_n]$, $\sum\deg f_i<n$. The number of common zeros in $\mathbb{F}_q^n$ is divisible by $p$.
Consequence: if all $f_i$ are homogeneous with no constant term, there is a non-trivial common zero (see also \ref{subsec:chevalley-warning-linalg}).

\subsection{Combinatorial Nullstellensatz (Alon)}
\label{subsec:combinatorial-nullstellensatz-poly}

$F$ is a field, $S_i\subset F$ finite, $f\in F[x_1,\dots,x_n]$, $\deg f\le\sum t_i$ with $t_i<|S_i|$. If the coefficient of $f$ at $\prod x_i^{t_i}$ is non-zero, then $f$ takes a non-zero value on $S_1\times\dots\times S_n$ (see also \ref{subsec:combinatorial-nullstellensatz}).

\subsection{Erd\H{o}s--Ginzburg--Ziv}
\label{subsec:egz-poly}

Among any $2n-1$ integers there are $n$ whose sum is divisible by $n$ (see also \ref{subsec:egz-combinatorics}).

\medskip

\section{Field extensions}
\label{sec:field-extensions}

\subsection{Minimal polynomial}
\label{subsec:minimal-polynomial-field}

For algebraic $\alpha\in L/K$ there is a unique monic $m_\alpha\in K[x]$ of smallest degree with $m_\alpha(\alpha)=0$. Properties: $m_\alpha$ is irreducible over $K$; $K[\alpha]=K(\alpha)$, $[K(\alpha):K]=\deg m_\alpha$; $\{1,\alpha,\dots,\alpha^{n-1}\}$ is a basis.

\subsection{Conjugates and symmetric functions}
\label{subsec:conjugates}

Let $m_\alpha=\prod_{i=1}^n(x-\alpha_i)$ (separable). For any $f\in K[x]$: $\sum f(\alpha_i)\in K$, $\prod f(\alpha_i)\in K$.

\subsection{Algebraic integers}
\label{subsec:algebraic-integers}

If an algebraic integer is rational, it is an integer.

\subsection{Trace and norm}
\label{subsec:trace-norm}

For a separable $L/K$ of degree $n$ with embeddings $\sigma_i$:
$$\operatorname{Tr}_{L/K}(\alpha)=\sum\sigma_i(\alpha),\quad \operatorname{N}_{L/K}(\alpha)=\prod\sigma_i(\alpha).$$
For $L=K(\alpha)$ with $m_\alpha=x^n+c_{n-1}x^{n-1}+\dots+c_0$: $\operatorname{Tr}(\alpha)=-c_{n-1}$, $\operatorname{N}(\alpha)=(-1)^n c_0$.
Transitivity: $\operatorname{Tr}_{M/K}=\operatorname{Tr}_{L/K}\circ\operatorname{Tr}_{M/L}$.

\subsection{Structure of $\mathbb{F}_q$}
\label{subsec:structure-fq}

For $q=p^n$ there exists a unique field $\mathbb{F}_q$. $\mathbb{F}_q=\{a\in\overline{\mathbb{F}}_p:a^q=a\}$ is the splitting field of $x^q-x$. $\mathbb{F}_q^\times$ is cyclic.

\subsection{Cyclotomic polynomial $\Phi_n$}
\label{subsec:cyclotomic}

$$\Phi_n(x)=\prod_{\substack{1\le k\le n\\\gcd(k,n)=1}}(x-\zeta_n^k)\in\mathbb{Z}[x],\quad\deg\Phi_n=\varphi(n).$$
$\Phi_n$ is irreducible over $\mathbb{Q}$. $x^n-1=\prod_{d\mid n}\Phi_d(x)$, $\Phi_n(x)=\prod_{d\mid n}(x^d-1)^{\mu(n/d)}$.

\subsection{Tower formula}
\label{subsec:tower-formula}

$[L:K]=[L:F]\cdot[F:K]$.

\subsection{Primitive element}
\label{subsec:primitive-element}

A finite separable $L/K$ is of the form $L=K(\theta)$. For $K(\alpha,\beta)$ with infinite $K$, $\theta=\alpha+c\beta$ works for almost all $c$.

\subsection{Kronecker on roots of unity}
\label{subsec:kronecker-roots}

If $\alpha\ne0$ is an algebraic integer and all its conjugates satisfy $|\alpha_i|\le1$, then $\alpha$ is a root of unity.

\medskip

\section{Factorization and irreducibility}
\label{sec:irreducibility}

\subsection{Gauss's lemma}
\label{subsec:gauss-lemma}

The product of primitive polynomials in $\mathbb{Z}[x]$ is primitive. $f$ is irreducible over $\mathbb{Z}$ $\Leftrightarrow$ irreducible over $\mathbb{Q}$ (up to a constant factor).

\subsection{Eisenstein's criterion}
\label{subsec:eisenstein}

$f=a_nx^n+\dots+a_0\in\mathbb{Z}[x]$, $p$ prime. If $p\nmid a_n$, $p\mid a_i$ ($0\le i\le n-1$), $p^2\nmid a_0$, then $f$ is irreducible over $\mathbb{Q}$.
\textit{Tricks:} shift $x\mapsto x+a$; reversal $\hat f(x)=x^n f(1/x)$.

\subsection{Generalized Eisenstein}
\label{subsec:eisenstein-generalized}

If $p\mid a_0,\dots,a_k$, $p\nmid a_{k+1}$, $p^2\nmid a_0$, then every non-constant factor of $f$ has degree $>k$.

\subsection{Reduction modulo $p$}
\label{subsec:reduction-mod-p}

$f\in\mathbb{Z}[x]$, $p\nmid a_n$. If $\bar f$ is irreducible in $\mathbb{F}_p[x]$, then $f$ is irreducible in $\mathbb{Z}[x]$. The converse is false ($x^4+1$ is a counterexample).

\subsection{Root estimate (simple constant term)}
\label{subsec:root-estimate-irred}

$f=a_nx^n+\dots+a_0\in\mathbb{Z}[x]$, $|a_0|$ is prime, $|a_0|>|a_1|+\dots+|a_n|$. Then $f$ is irreducible.

\subsection{Irreducibility of $\Phi_p$ via shift}
\label{subsec:phi-p-eisenstein}

$\Phi_p(x+1)=\frac{(x+1)^p-1}{x}=x^{p-1}+\binom{p}{1}x^{p-2}+\dots+\binom{p}{p-1}$. Apply Eisenstein with $p$.

\subsection{Newton polygon}
\label{subsec:newton-polygon}

$f=\sum a_i x^i\in\mathbb{Z}_p[x]$, plot points $(i,v_p(a_i))$, take lower convex hull. If the slopes are $\lambda_1,\dots,\lambda_k$ with multiplicities $m_1,\dots,m_k$, then the valuations of the roots are $-\lambda_i$ with multiplicity $m_i$. Irreducibility criterion: a single segment of slope $m/n$ in lowest terms $\Rightarrow$ $f$ remains irreducible over any extension of degree $<n$. If $v_p(a_n)=0$, $v_p(a_0)=1$, and every other $v_p(a_i)\ge1$, then $f$ is irreducible (Eisenstein via Newton).

\subsection{Hensel lifting}
\label{subsec:hensel-polynomial}

$f\in\mathbb{Z}[x]$, $f(a)\equiv0\pmod p$, $f'(a)\not\equiv0\pmod p$. Then for any $k\ge1$ there exists a unique $\tilde a\pmod{p^k}$ with $\tilde a\equiv a\pmod p$ and $f(\tilde a)\equiv0\pmod{p^k}$ (see also \ref{subsec:hensel}).
